\section*{Chapter \ref{ch:dv:rough-volatility-and-forward-variance-models} --- Rough Volatility and Forward-Variance Models}

\begin{solution}{exo:dv:rough-volatility-and-forward-variance-models:1}
Total variances are $0.04$ and $2\times0.0484=0.0968$; the \omterm{def:dv:rough-volatility-and-forward-variance-models:fwdvar}{forward variance} is $0.0568$ a year, a forward
volatility of 23.8\%.
\end{solution}

\begin{solution}{exo:dv:rough-volatility-and-forward-variance-models:2}
The strip weights each out-of-the-money option by $1/K^2$, so low strikes count more than high ones. With a
downside skew the low-strike puts are priced at higher volatilities than the at-the-money option, and the
weighted average of variance is above the at-the-money variance.
\end{solution}

\begin{solution}{exo:dv:rough-volatility-and-forward-variance-models:3}
$m(2,\Delta)\propto\Delta^{2H}$, so $2^{2H}=1.15$ and $H=\ln1.15/(2\ln2)=0.10$.
\end{solution}

\begin{solution}{exo:dv:rough-volatility-and-forward-variance-models:4}
$Y_t$ is Gaussian with mean zero and variance $2H\int_0^t(t-s)^{2H-1}ds=t^{2H}$, so
$\E[e^{\eta Y_t}]=e^{\eta^2t^{2H}/2}$ and $\E[v_t]=\xi_0(t)$. For $u\ge t$,
$\xi_t(u)=\E_t[v_u]$: split $Y_u$ into $\sqrt{2H}\int_0^t(u-s)^{H-1/2}dW_s$, known at $t$, and an independent
remainder of variance $(u-t)^{2H}$; taking the expectation of the remainder's exponential gives the formula of
the text. Its only $t$-dependence through $W$ is $\eta\sqrt{2H}\int_0^t(u-s)^{H-1/2}dW_s$, whose differential is
$\eta\sqrt{2H}(u-t)^{H-1/2}dW_t$: that is the volatility of $\xi_t(u)$.
\end{solution}

\begin{solution}{exo:dv:rough-volatility-and-forward-variance-models:5}
The future is $\E[\mathrm{VIX}]$ and the forward volatility $\sqrt{\E[\mathrm{VIX}^2]}$; by Jensen's inequality
$\E[\mathrm{VIX}]\le\sqrt{\E[\mathrm{VIX}^2]}$, with a gap that grows with the variance of the VIX. A larger
$\eta$ makes the VIX more volatile and widens the gap: with $\eta=2.5$ the future is 17.88.
\end{solution}

\begin{solution}{exo:dv:rough-volatility-and-forward-variance-models:6}
In the calm state $R_1=0$ and $R_2=\sigma^2$, so $\sigma^*=\beta_0+\beta_2\sigma^*$ and $\sigma^*=0.04/0.35=11.4\%$.
After a $+4\%$ day, $\beta_1R_1$ is negative and larger in size than the rise of $\beta_2\sqrt{R_2}$, so
volatility falls; $R_1$ decays at rate 25 and $R_2$ at 15, so the activity term outlasts the trend term and
volatility rises above $\sigma^*$, to 12.4\%, before both fade.
\end{solution}

\begin{solution}{exo:dv:rough-volatility-and-forward-variance-models:7}
Over five seeds: 0.45 to 0.50 without noise, the diffusion's $\frac12$ lowered a little by mean reversion; 0.20
to 0.23 with noise. The rough path of the same code reads 0.10 to 0.12.
\end{solution}

\begin{solution}{exo:dv:rough-volatility-and-forward-variance-models:8}
Daily realised volatility is measured with an error of about 0.08 in log, independent from day to day. The error
adds a constant to the structure function at every lag and flattens its slope: a diffusion with $H=\frac12$
read through it gives about 0.17 to 0.23. The estimate refutes nothing until the noise is modelled, for instance by subtracting its contribution, twice
its variance, from $m(2,\Delta)$ before the regression; the evidence
from the skew term structure, where no measurement noise enters, is the stronger test.
\end{solution}

\begin{solution}{pb:dv:rough-volatility-and-forward-variance-models:1}
\textbf{1.} 22.1\% and 19.2\%.
\textbf{2.} 24.0\%.
\textbf{3.} 25.8\%.
\textbf{4.} 24.9\%.
\textbf{5.} Yes: total variance rises at every pillar, and \texttt{negative\_intervals()} is empty.
\textbf{6.} 4.13, 2.30 and 0.85.
\textbf{7.} $\eta=2.51$, $\kappa=1.08$.
\textbf{8.} 3.1 (7.76 against 2.50).
\textbf{9.} The exponential is at most $\eta$ at zero horizon; the power law is infinite there. Two points fix
$\eta$ and $\kappa$, and the curves then differ everywhere else.
\textbf{10.} It adds a fast-decaying term that lifts the short end; a sum of exponentials can imitate the power
law over a chosen range of horizons, never down to zero.
\textbf{11.} $-1.44$ and $-0.25$.
\textbf{12.} $-0.70$ and $-0.25$.
\textbf{13.} $-0.44$ for both.
\textbf{14.} $-0.09$: nearly flat.
\textbf{15.} $1.44/0.25=5.7$, and a power law with exponent $-0.44$ predicts $52^{0.44}=5.7$ between one week and one
year.
\textbf{16.} $H=\frac12-0.44=0.06$.
\textbf{17.} The rule applied to a model with a known $H=0.1$ returns 0.06 as well: over this range of expiries it
reads low by about 0.04, so the surface is consistent with $H$ near 0.1.
\textbf{18.} Time series give 0.075 to 0.158 across 21 indices, 0.124 for the S\&P~500 in the same table: consistent
with 0.1.
\textbf{19.} Exponent $-0.44$; $H=0.06$ by the short-expiry rule, about 0.1 once the rule is calibrated on rough
Bergomi.
\textbf{20.} Yes for short-dated index options, whose skew Heston cannot produce. Costs: Monte Carlo only (no PDE,
slower and noisier \omterm{def:dv:greeks-and-the-hedging-pnl:greeks}{Greeks}), and the VIX smile the model implies must be checked against the market's before its
volatility-of-volatility products are trusted.
\end{solution}

\begin{solution}{iq:dv:rough-volatility-and-forward-variance-models:1}
$\xi_t(u)=\E^{\mathbb Q}_t[v_u]$, the expected instantaneous variance at a future date. The expected integrated variance
to $T$ is minus twice the log contract's price, which a strip of out-of-the-money options replicates; that is the
variance-swap strike $T\sigma^2_{\mathrm{VS}}(T)$, and its derivative in $T$ is $\xi_0(T)$.
\iqlookfor{the definition, the log-contract strip, and the derivative in expiry.}
\end{solution}

\begin{solution}{iq:dv:rough-volatility-and-forward-variance-models:2}
The logarithm of volatility behaves at short time scales like a fractional Brownian motion with $H$ near 0.1, far
rougher than a diffusion's $\frac12$. Evidence: the scaling of the structure function of realised volatility across
many indices and moments, and the power-law at-the-money skew $\psi\sim T^{H-1/2}$ in option prices. Caveat: the
time-series estimate is sensitive to measurement noise.
\iqlookfor{both kinds of evidence, and the caveat.}
\end{solution}

\begin{solution}{iq:dv:rough-volatility-and-forward-variance-models:3}
In a diffusive \omterm{def:dv:stochastic-volatility:sv}{stochastic volatility model} the short skew comes from the volatility of variance over a short time,
which is bounded, so the skew tends to a constant as $T\to0$. A \omterm{def:dv:rough-volatility-and-forward-variance-models:rough}{rough volatility} has a volatility of \omterm{def:dv:rough-volatility-and-forward-variance-models:fwdvar}{forward
variance} that explodes like $\tau^{H-1/2}$ at short horizons, producing a skew $\sim T^{H-1/2}$ that keeps
growing, like the market's.
\iqlookfor{the short-expiry limit, and the kernel.}
\end{solution}

\begin{solution}{iq:dv:rough-volatility-and-forward-variance-models:4}
The model is not Markov, so Monte Carlo. Exact: the Volterra process and its Brownian motion are jointly Gaussian on
the grid, so Cholesky of the $2n\times2n$ covariance, $O(n^3)$ once and $O(n^2)$ per path. Faster: the hybrid scheme,
exact near the singularity and a fast-Fourier convolution elsewhere, $O(n\log n)$. Price with the \omterm{prop:dv:stochastic-volatility:mixing}{mixing formula}
(conditioning on the variance driver) to cut noise; check the martingale property.
\iqlookfor{non-Markov, the Gaussian structure, the costs, variance reduction.}
\end{solution}

\begin{solution}{iq:dv:rough-volatility-and-forward-variance-models:5}
Fitting the index smile and the VIX smile with one model. The index smile constrains the skew and the volatility of
short-dated variance; the VIX smile constrains the law of that variance directly, and simple models cannot do both
(Guyon's conjecture; answered by the quadratic rough \omterm{def:dv:stochastic-volatility:heston}{Heston model}, by path-dependent volatility and by Guyon's exact
construction). A desk cares because VIX products and index products sensitive to \omterm{def:dv:stochastic-volatility:sv}{volatility of volatility}
(cliquets, options on variance) are priced on exactly the dynamics a one-sided fit leaves untested.
\iqlookfor{what each smile constrains, and the products at risk.}
\end{solution}

\begin{solution}{iq:dv:rough-volatility-and-forward-variance-models:6}
They agree on vanillas of the quoted expiries and differ on everything that depends on the volatility of \omterm{def:dv:rough-volatility-and-forward-variance-models:fwdvar}{forward
variance} at other horizons: short-dated options below the quoted expiries, forward-start options and cliquets
(forward skew), options on variance and on the VIX. Decide by testing each against the instruments that measure
those dynamics (weekly options, VIX futures and options, forward-start quotes) and by comparing the \omterm{def:dv:greeks-and-the-hedging-pnl:hpnl}{hedging P\&L}
each model would have produced; hold a model reserve for the difference on the products that remain.
\iqlookfor{the product list, a test against other markets, and a reserve.}
\end{solution}
